\subsection{Candidate space}

The confirmatory bank evaluated every resolved candidate in each profile ---
\ConfirmResourceConstrainedCandidates{}, \ConfirmIntermediateCandidates{},
and \ConfirmHighCapacityCandidates{} in the resource-constrained,
intermediate-capacity, and high-capacity profiles --- on
\ConfirmResourceConstrainedScenarios{}, \ConfirmIntermediateScenarios{}, and
\ConfirmHighCapacityScenarios{} fresh paired scenarios each, for
\ConfirmTotalExecutions{} executions in total.

The candidate counts fall sharply with profile strength, and that is the
first substantive consequence of separating capacity from controls. A
high-capacity environment already owns least-privilege segmentation and
isolated backups, so those are not decisions it can make. What remains is
patching, detection, isolation, and identity controls:
\ConfirmHighCapacityCandidates{} distinct configurations rather than the
\EnumeratedPortfolios{} the enumeration nominally contains. Under the
previous version's precedence rule the same profile appeared to face far
more choices, many of them cheaper than doing nothing, because a candidate
could sell back the segmentation and backup architecture the profile already
had.

Exactly one candidate per profile costs nothing: the profile's own untouched
posture. Its mean weighted service-hours lost are
\ConfirmResourceConstrainedFreeBaselineMeanLoss{},
\ConfirmIntermediateFreeBaselineMeanLoss{}, and
\ConfirmHighCapacityFreeBaselineMeanLoss{} respectively, and these are the
reference points against which every purchase should be read.

\subsection{Full-space versus finalist-only frontiers}

Over the full space, \ConfirmFullSpaceFrontier{} candidates are
non-dominated: \ConfirmResourceConstrainedNonDominated{},
\ConfirmIntermediateNonDominated{}, and \ConfirmHighCapacityNonDominated{} by
profile. Restricting attention to the \ConfirmFinalistLabelled{} candidates
that the frozen discovery rule labels as finalists, and recomputing dominance
within that subset alone, gives a frontier of \ConfirmFinalistOnlyFrontier{}.

The two views disagree in both directions, and the disagreement is the point
(Figure~\ref{fig:fullspace}).

\textbf{\ConfirmFinalistOnlyArtifacts{} candidates are efficient only in the
restricted view.} They appear non-dominated among the finalists and are
dominated once the full space is present. In each case the dominating
candidate is instructive rather than exotic. In the high-capacity profile the
displaced candidate is dominated by the portfolio identical to it except for
dropping identity controls: within this model, in that environment, identity
controls cost four points and buy nothing detectable. Two
intermediate-capacity candidates are displaced by portfolios that retain the
profile's own periodic backup architecture instead of upgrading it.

\textbf{\ConfirmNonFinalistsOnFrontier{} non-dominated candidates were never
labelled as finalists at all.} A design that evaluated only the finalist set
could not have shown them to a reader under any circumstances, and could not
have reported that they existed. Of the \ConfirmFinalistLabelled{} labelled
finalists, \ConfirmFinalistsSurviving{} remain non-dominated over the full
space and \ConfirmFinalistsDisplaced{} do not.

We report these magnitudes as modest, because they are. The discovery rule is
not badly wrong: most of what it selects survives, and most of the frontier
it would have shown is genuinely efficient. What it cannot do is tell a
reader which of those two statements is true, because a restricted comparison
contains no evidence about the candidates it excluded. That is a structural
property of the design rather than a property of this dataset, and it is why
the count is worth having at all.

\begin{figure*}[t]
\centering
\includegraphics[width=\textwidth]{full_space_vs_finalist.pdf}
\caption{Candidate counts at each stage. The gap between the third and fourth
bars is an effect of the candidate set on the reported frontier, not an
effect of the data. Panel subtitles give the number of candidates that appear
efficient within the finalist subset but are dominated once every resolved
candidate is present.}
\label{fig:fullspace}
\end{figure*}

\subsection{Where the decision is live, and where it is not}

The three profiles do not merely differ in degree; they pose different kinds
of decision (Figure~\ref{fig:frontier}).

\textbf{High-capacity: nothing worth buying.} Every one of the
\ConfirmHighCapacityCandidates{} resolved candidates falls between
\ConfirmHighCapacityMeanLossMin{} and \ConfirmHighCapacityMeanLossMax{}
weighted service-hours. The entire purchasable improvement in this profile is
\ConfirmHighCapacityMaxPurchasableImprovement{} weighted service-hours, and
the best candidate within the \ConfirmHighCapacityBudget-point budget spends
\ConfirmHighCapacityBestAffordableCost{} points to reach
\ConfirmHighCapacityBestAffordableMeanLoss{} hours from a free baseline of
\ConfirmHighCapacityFreeBaselineMeanLoss{}. Against a Monte Carlo standard
error of \ConfirmHighCapacityWorstMeanLossStandardError{} hours, that
improvement sits at the edge of what \ConfirmHighCapacityScenarios{} paired
scenarios can resolve:
\ConfirmHighCapacityIndistinguishablePairs{} of
\ConfirmHighCapacityCandidatePairs{} candidate pairs
(\ConfirmHighCapacityIndistinguishableShare\%) are separated by less than one
combined standard error on every stochastic objective. The defensible reading
is that in this modeled environment no purchasable control makes a detectable
difference to disruption, and the frontier is driven by cost and burden
alone.

\textbf{Intermediate-capacity: the decision that matters.} The free baseline
loses \ConfirmIntermediateFreeBaselineMeanLoss{} hours and the best candidate
within the \ConfirmIntermediateBudget-point budget loses
\ConfirmIntermediateBestAffordableMeanLoss{} at a cost of
\ConfirmIntermediateBestAffordableCost{} points --- a modeled improvement of
more than an order of magnitude, with sustained-outage probability at
\ConfirmIntermediateBestAffordableOutage\%. This is the profile in which the
six objectives genuinely trade against one another and in which a preference
statement changes the answer.

\textbf{Resource-constrained: the budget does not reach.} The free baseline
loses \ConfirmResourceConstrainedFreeBaselineMeanLoss{} hours. The best
candidate anywhere in the space reaches
\ConfirmResourceConstrainedMeanLossMin{}, but the best candidate that fits
the declared \ConfirmResourceConstrainedBudget-point budget reaches only
\ConfirmResourceConstrainedBestAffordableMeanLoss{}, and its sustained-outage
probability is \ConfirmResourceConstrainedBestAffordableOutage\%. Within
budget that endpoint is close to immovable: the affordable options change
mean disruption substantially while leaving multi-service clinical outage
nearly certain. The honest statement is not that some affordable portfolio is
adequate. It is that in this modeled environment the declared budget does not
buy the outcome the endpoint measures.

\begin{figure*}[t]
\centering
\includegraphics[width=\textwidth]{confirmatory_frontier.pdf}
\caption{Cost against mean weighted service-hours lost, with Monte Carlo
standard errors. Lines connect full-space non-dominated candidates in this
two-objective projection; dominance is evaluated in all six. Note the
vertical scales: the high-capacity panel spans under one service-hour.}
\label{fig:frontier}
\end{figure*}

\subsection{Uncertainty and resolution}

Worst-case Monte Carlo standard error on mean disruption is
\ConfirmResourceConstrainedWorstMeanLossStandardError{},
\ConfirmIntermediateWorstMeanLossStandardError{}, and
\ConfirmHighCapacityWorstMeanLossStandardError{} weighted service-hours in
the three profiles.

The share of candidate pairs the scenario bank cannot reliably order is
\ConfirmResourceConstrainedIndistinguishableShare\%,
\ConfirmIntermediateIndistinguishableShare\%, and
\ConfirmHighCapacityIndistinguishableShare\%. The per-profile scenario
allocation was chosen precisely to attack this. At the
\DiscoveryHighCapacityScenarios{} scenarios used in discovery, every one of
the \DiscoveryHighCapacityCandidatePairs{} high-capacity pairs was
indistinguishable; raising that profile alone to
\ConfirmHighCapacityScenarios{} scenarios reduced it to
\ConfirmHighCapacityIndistinguishablePairs{}. A uniform allocation would have
spent scenarios where they were already ample and withheld them where they
were scarce.

Bootstrap Pareto inclusion is at least \StabilityThreshold{} for
\ConfirmResourceConstrainedStableAtEighty{},
\ConfirmIntermediateStableAtEighty{}, and
\ConfirmHighCapacityStableAtEighty{} candidates respectively.

\subsection{The sustained-outage endpoint at every $k$}

Table~\ref{tab:kladder} reports the endpoint at every clinical service count.
The values move by roughly three percentage points across the whole range of
$k$, for the reason given in the methods: within this model the clinical
estate is close to all-or-nothing (Figure~\ref{fig:bimodal}). A reader who
disagrees with the primary $k = \PrimaryK$ can substitute any other and reach
substantially the same conclusion. That is a statement about the model's
behaviour, not a defence of the choice.

\begin{table}[t]
\centering
\caption{Mean sustained-outage probability across candidates, at every
clinical service count $k$. The primary analysis uses $k = \PrimaryK$.}
\label{tab:kladder}
\small
\begin{tabular}{@{}lrrrr@{}}
\toprule
Profile & $k=1$ & $k=2$ & $k=3$ & $k=4$ \\
\midrule
Resource-constrained & \ConfirmResourceConstrainedOutageKOneMean\% & \ConfirmResourceConstrainedOutageKTwoMean\% & \ConfirmResourceConstrainedOutageKThreeMean\% & \ConfirmResourceConstrainedOutageKFourMean\% \\
Intermediate & \ConfirmIntermediateOutageKOneMean\% & \ConfirmIntermediateOutageKTwoMean\% & \ConfirmIntermediateOutageKThreeMean\% & \ConfirmIntermediateOutageKFourMean\% \\
High-capacity & \ConfirmHighCapacityOutageKOneMean\% & \ConfirmHighCapacityOutageKTwoMean\% & \ConfirmHighCapacityOutageKThreeMean\% & \ConfirmHighCapacityOutageKFourMean\% \\
\bottomrule
\end{tabular}
\end{table}

\begin{figure}[t]
\centering
\includegraphics[width=\columnwidth]{clinical_outage_bimodality.pdf}
\caption{Number of clinical services sustaining a qualifying outage across
the discovery bank. The near-binary shape is why the choice of $k$ is not
load-bearing --- and it is an artifact of binary per-step service
availability, not a property of real incidents.}
\label{fig:bimodal}
\end{figure}

\subsection{External validation}

Of \ValidationBenchmarks{} frozen external benchmarks, \ValidationScored{}
could be scored numerically. \ValidationNotAddressable{} are not addressable
because the model has no output of that construct at any parameter setting,
and \ValidationNotScored{} cannot be scored because the model's quantity and
the target's quantity are not the same thing. \ValidationFailed{} benchmarks
fail. None passes.

That ratio is the headline validation result and it bounds every conclusion
above. Three failures are worth stating individually.

\textbf{The restoration engine is anchored to neither reference.} Against the
one directly observed large-scale technical hospital outage the model is far
too slow: its median clinical outage is
\ValidationOutageModelMedianHours{} modeled hours against an observed median
of \ValidationOutageTargetMedianHours{}, and
\ValidationOutageModelSixHourShare\% of disrupted trials recover within six
hours against an observed \ValidationOutageTargetSixHourShare\%
\citep{tully2025}. Against the ransomware incident record the same engine is
far too fast, completing inside a \HorizonHours-hour horizon where the public
duration distribution has most of its mass beyond a week
\citep{neprash2022}. A model simultaneously too slow for a technical outage
and too fast for a ransomware incident is anchored to neither, and its
non-recovery objective should be read as a within-model ordering device
rather than a duration estimate.

\textbf{Multi-facility breadth fails outright.} The simulator runs one
facility, so it assigns probability zero to roughly a quarter of publicly
reported hospital-linked events --- those that affected more than one
Medicare-certified hospital. This is a total discrepancy on a directly
comparable target.

\textbf{Regional spillover cannot be addressed at all.} The model has no
representation of a neighboring hospital, patient transfer, or diversion, so
the documented burden at adjacent emergency departments
\citep{dameff2023,abouk2024} has no counterpart output to compare against.

The one benchmark the model is not obviously wrong on is the share of
incidents contained before clinical service loss, where the model's spread
across profiles brackets the surveyed figure. We report that as consistency
rather than a pass: a self-reported vendor survey spanning many countries and
organization sizes is not a U.S. hospital denominator, and bracketing a point
estimate with a range that wide is weak evidence.

\begin{figure}[t]
\centering
\includegraphics[width=\columnwidth]{external_validation_scorecard.pdf}
\caption{Verdicts against the frozen benchmark registry. Targets and criteria
were fixed before any comparison was run, and failures are reported rather
than corrected.}
\label{fig:validation}
\end{figure}

\subsection{Structural sensitivity}

Assumption S1 --- that restoration begins only once the estate is contained
--- was varied on the candidates carrying the conclusions above. Of
\StructuralCandidatesTested{} candidates evaluated under both structures,
\StructuralPreserved{} keep their frontier membership;
\StructuralOnlyGated{} are non-dominated only under the primary structure and
\StructuralOnlyParallel{} only under the alternative. The worst per-profile
agreement is \StructuralWorstAgreement\%.

Structural uncertainty is not expressible as an interval, and this comparison
is the substitute: the conclusions above are conditional on S1, and the
magnitude of that conditionality is measured rather than assumed away.
